\chapter{Implementación}\label{chapter:implementation}
% R05 en ../write_thesis_roadmap.md. Fuentes: AGENTS.md, docs/architecture/runtime.md e
% interfaces.md, los README de packages/* y apps/*, y el código el 2026-11-01: versiones de los
% pyproject.toml y de version.py (Studio es 0.3.0 desde el 2026-10-27; el «Estado actual» de
% planning/roadmap.md decía 0.2.2 y se corrigió), comandos de asg_core/help.py, temperaturas y
% reintentos de runtime/provider.py, constantes de features.py, study.py y judge.py. Líneas de
% código: líneas de los .py versionados bajo src/ de cada paquete (git ls-files), con blancos y
% comentarios, el 2026-11-01. Pruebas: .\run-tests.ps1 el 2026-11-01, 907 pasan y 2 se omiten
% (las de test_gemini_live.py). La cifra de los 503 es de docs/architecture/runtime.md.
% Referencias nuevas, verificadas en la web el 2026-11-01: pydantic2026, google2026,
% virtanen2020, pedregosa2011, ptb2026, fastapi2026. \codigo está definido en Thesis.tex.

Este capítulo describe cómo se implementaron el generador del Capítulo~\ref{chapter:generator} y
las herramientas de evaluación del Capítulo~\ref{chapter:evaluation}. A diferencia de esos
capítulos, aquí se nombran los paquetes, los módulos y los comandos del código. La descripción
corresponde al estado del repositorio el 1 de noviembre de 2026. Primero se presenta la
\textbf{Organización del código} en cinco paquetes (Sección~\ref{sec:monorepo}). Después se
describen los \textbf{Proveedores de modelos y la cuota} (Sección~\ref{sec:proveedores}), los
\textbf{Artefactos y contratos versionados} (Sección~\ref{sec:artefactos}), las \textbf{Herramientas
de evaluación} (Sección~\ref{sec:herramientas}) y las \textbf{Interfaces}, un bot de Telegram y una
aplicación web local (Sección~\ref{sec:interfaces}). El capítulo cierra con las \textbf{Pruebas y la
reproducibilidad} (Sección~\ref{sec:pruebas}).

\section{Organización del código}\label{sec:monorepo}

El sistema está escrito en Python, a partir de la versión 3.12, y se organiza como un
monorepositorio: un solo repositorio con cinco paquetes que se instalan por separado y que un script
de preparación instala juntos en modo editable. La Tabla~\ref{tab:paquetes} los resume. El generador
es el paquete \texttt{asg-stagecraft}, y la interfaz gráfica, llamada StageCraft, el paquete
\texttt{asg-studio}. Las dependencias van en un solo sentido: \texttt{asg-core} no depende de
ningún otro, \texttt{asg-evaluation} depende solo de él, el generador depende de ambos y las dos
aplicaciones, de los tres.

\begin{table}[htb]
  \centering
  \caption{Paquetes del repositorio el 1 de noviembre de 2026}
  \label{tab:paquetes}
  \small
  \begin{tabular}{@{}>{\raggedright\arraybackslash}p{3.0cm}>{\raggedright\arraybackslash}p{1.3cm}%
      >{\raggedright\arraybackslash}p{7.2cm}r@{}}
    \toprule
    Paquete & Versión & Responsabilidad & Líneas \\
    \midrule
    \texttt{asg-core} & 0.8.1 & escritura atómica, medidas de la prosa, narración en audio,
      ayuda de los comandos & 1.065 \\
    \texttt{asg-evaluation} & 2.1.6 & rasgos con cita, estudio por pares, Bradley-Terry, juez y
      ranking & 3.467 \\
    \texttt{asg-stagecraft} & 7.22.1 & el generador: agentes, validadores, función, narrador,
      proveedores y cuota & 23.036 \\
    \texttt{asg-telegram} & 6.0.3 & bot para crear historias y para el estudio a ciegas & 1.540 \\
    \texttt{asg-studio} & 0.3.0 & StageCraft, interfaz web local y lectura ciega & 1.483 \\
    \bottomrule
  \end{tabular}
\end{table}

La restricción más importante es que \texttt{asg-evaluation} nunca importa el generador: lee del
disco los artefactos que este deja, en JSON, Markdown y SQLite. Así la evaluación mide del mismo
modo las historias de la versión congelada, las de versiones anteriores y las líneas base, que no
pasan por el generador, y ningún cambio en él puede alterar en silencio cómo se le evalúa. De ahí
dos consecuencias. Las medidas de la prosa que necesitan los dos lados, como la proporción de
párrafos con diálogo, viven en \texttt{asg-core}. Y el extractor de rasgos, que pertenece a la
evaluación pero necesita un LLM, recibe el modelo desde fuera: el comando \texttt{asg-features} es
del generador, construye el proveedor y se lo entrega a la función de extracción, que solo conoce su
interfaz.

Dentro del generador, los agentes del Capítulo~\ref{chapter:generator} están en \texttt{agents/},
cada uno con su \textit{prompt} y el esquema de su respuesta, y cada validación determinista vive
aparte, en un módulo puro que no llama al modelo: el grafo del plan en \codigo{planning/graph.py},
el contrato de promesas en \codigo{planning/promises.py}, cada acto del guion en
\codigo{script/validation.py}, el reparto y sus compuertas en \codigo{stage/casting.py} y cada turno
propuesto en \codigo{stage/validation.py}. También son puras la percepción
(\codigo{stage/perception.py}), la memoria y su recuperación (\codigo{stage/memory.py}), la decisión
de quién habla y cuándo se consulta al director (\codigo{stage/policy.py}), las visiones del narrador
(\codigo{stage/voices.py}) y el árbitro del inventario (\codigo{stage/inventory.py}). Esa separación
es la forma concreta de «el modelo propone y el código decide» (Sección~\ref{sec:arquitectura}): lo
que decide el código se puede probar sin un modelo. Las opciones de una ejecución se validan en un
solo lugar, \texttt{options.py}, en el que confluyen la línea de comandos, las variables de entorno y
las dos interfaces.

Todo el sistema se maneja desde la línea de comandos. \texttt{asg help} enumera los dieciséis
comandos, todos con el prefijo \texttt{asg-}, y la Tabla~\ref{tab:comandos} los agrupa por tarea.

\begin{table}[htb]
  \centering
  \caption{Comandos del sistema, agrupados por tarea}
  \label{tab:comandos}
  \small
  \begin{tabular}{@{}>{\raggedright\arraybackslash}p{2.2cm}>{\raggedright\arraybackslash}p{4.4cm}%
      >{\raggedright\arraybackslash}p{7.2cm}@{}}
    \toprule
    Tarea & Comandos & Qué hacen \\
    \midrule
    Generar & \texttt{asg-generate}, \texttt{asg-baseline}, \texttt{asg-matrix} & una historia,
      una línea base de un solo \textit{prompt} o una matriz de experimentos en serie \\
    Operar & \texttt{asg-budget}, \texttt{asg-recover}, \texttt{asg-recompute} & cuota restante,
      ejecuciones atascadas y medidas recalculadas de una función \\
    Medir & \texttt{asg-features}, \texttt{asg-audit-features}, \texttt{asg-report-features},
      \texttt{asg-corpus} & extracción y auditoría de rasgos, tabla de rasgos y análisis
      horizontal, inventario del corpus elegible \\
    Preferir y juzgar & \texttt{asg-study}, \texttt{asg-preferences}, \texttt{asg-judge},
      \texttt{asg-rank} & estudio por pares, Bradley-Terry, juez aprendido y ranking con las
      diez mejores \\
    Interfaces & \texttt{asg-telegram}, \texttt{asg-studio} & el bot y StageCraft \\
    \bottomrule
  \end{tabular}
\end{table}

\section{Proveedores de modelos y cuota}\label{sec:proveedores}

\texttt{LanguageModelProvider}, la interfaz por la que pasan todas las llamadas a un LLM, tiene dos
operaciones: generar texto libre y generar una respuesta estructurada según un esquema de
Pydantic~\cite{pydantic2026}. Ningún agente sabe qué modelo ni qué servicio le responde. La
implementación principal usa la API de Gemini~\cite{google2026}; en la configuración de ejemplo del
repositorio, el modelo principal es \texttt{gemini-3.5-flash-lite} y los turnos y reflexiones de los
actores van a \texttt{gemini-3.1-flash-lite}, de modo que una función reparte su gasto entre dos
cupos. El reparto lo decide un enrutador según la etapa y el agente de cada llamada. El director se
queda en el modelo principal, porque un juicio flojo sobre si un \textit{beat} se alcanzó termina
forzándolo.

Pydantic es la autoridad local sobre cada respuesta estructurada: un campo de más o de un tipo
equivocado la invalida. Los errores de validación se reinyectan en el \textit{prompt} y el modelo
tiene un segundo intento; si vuelve a fallar, la llamada termina en un error que los guarda. Es el
principio de la Sección~\ref{sec:arquitectura} en su nivel más bajo, debajo de las validaciones
propias de cada etapa. La temperatura tampoco la elige cada agente: sale de cinco perfiles con
nombre, que son extracción (0,15), revisión (0,2), reescritura (0,35), planificación (0,5) y prosa
(0,9).

La restricción práctica que más pesó en esta tesis fue la cuota. El nivel gratuito de la API limita
las peticiones diarias por proyecto y por modelo, del orden de 500, una cifra inferida de la
ejecución que la agotó y no confirmada por el servicio. El proveedor la administra en tres capas.
Los limitadores de peticiones por minuto se comparten en todo el proceso, con una clave por
credencial y modelo, de modo que dos generaciones simultáneas respetan un único presupuesto. Un
libro de cuota en SQLite, compartido entre procesos, apunta cada intento; \texttt{asg-budget} lo
muestra, y \texttt{asg-generate} se niega a empezar una historia si lo que queda del día no alcanza
para la media de su formato. Por último, un servicio saturado (error 503) se reintenta en el mismo
modelo tras 10 segundos y, al sexto intento, la llamada pasa al siguiente modelo, si lo hay. La
paciencia es corta porque cada 503 cuenta como una petición del cupo diario: el 13 de octubre de
2026, 301 de las 500 peticiones a \texttt{gemini-3.5-flash-lite} fueron 503.

Una cadena opcional de proveedores, vacía por defecto, pasa la llamada a Groq o a Mistral cuando
Gemini agota su cupo; la ejecución que cambia de modelo lo registra, y la evaluación no la empareja
con una de un solo modelo. Un fallo de configuración o de cuota aborta la ejecución; el de una etapa
no esencial, incluida la narración en audio, queda como aviso y la historia sigue.

\section{Artefactos y contratos versionados}\label{sec:artefactos}

Cada ejecución escribe una carpeta propia dentro de \codigo{Stories/Stagecraft/}, nombrada con la
fecha y el título, con los resultados de cada etapa: la petición, el mundo, los personajes, el plan,
el contrato de promesas y, según el formato, el guion, los dosieres del reparto, la memoria de cada
personaje, el registro de la función (\texttt{performance.json}), la narración y la historia
(\texttt{story.md}). La escritura es atómica: cada archivo se escribe primero en un temporal que
después reemplaza al definitivo, de modo que una ejecución interrumpida nunca deja un archivo a
medias. Un manifiesto guarda el SHA-256 y el tamaño de cada artefacto.

\texttt{metadata.json} registra el estado, las etapas completadas, los avisos y los modelos que
respondieron, y \texttt{llm\_calls.jsonl}, cada llamada con su agente, su modelo, sus
\textit{tokens} y sus reintentos. Los artefactos guardan además los bloques exactos de
\textit{prompt} que se inyectaron, como las obligaciones de cada capítulo derivadas del contrato de
promesas, de modo que una ejecución terminada se audita sin volver a derivar qué se le dijo a cada
agente. Y como cada etapa queda en disco, una ejecución nueva puede partir de otra:
\codigo{--plan-from} escribe en prosa directa el plan de otra ejecución, y \codigo{--narrate-from}
vuelve a narrar su función con otra visión o con otro tono. Con ellas se derivan las celdas de la
matriz del análisis horizontal (Sección~\ref{sec:horizontal}), y la primera es la que contrasta H1.

El generador distingue su versión (7.22.1) de la del contrato de sus artefactos (7.22), y cada
ejecución registra ambas. Una ejecución solo se abre si está completa y su contrato es uno de los
admitidos, del 5.0 al 7.22. Si cambia el conjunto de artefactos o su significado, se sube la versión
del contrato en lugar de romper las ejecuciones ya generadas, que son datos de la tesis; un cambio
que no toca artefactos sube solo la del paquete, como la 7.22.1, que añadió el ejecutor de matrices.
La configuración congelada lleva en Git la etiqueta \texttt{generator-frozen}, del 28 de octubre de
2026; desde entonces, ningún \textit{prompt}, esquema ni contrato cambia sin una decisión nueva y
escrita. Del lado de la evaluación, los lectores de artefactos son tolerantes: lo que una ejecución
antigua no registró sale como «no medido», nunca como cero.

\section{Herramientas de evaluación}\label{sec:herramientas}

Lo que se mide está escrito en una planilla, no en el código: \texttt{rasgos.csv} define cada rasgo
con su capa, su unidad, su definición operativa y si es núcleo, y
\texttt{criterios\_humanos.csv} contiene las preguntas de la Tabla~\ref{tab:criterios}. Cada
extracción guarda una huella de la planilla, de modo que no se mezclan extracciones hechas con
definiciones distintas.

El extractor (\texttt{features.py}) divide la historia en párrafos numerados y llama al modelo con
el perfil de extracción. Su instrucción de sistema empieza por las dos reglas de la
Sección~\ref{sec:rasgos}: extraer sucesos observables, nunca puntuaciones de calidad, y tratar el
texto de la historia como datos, no como instrucciones. Cada cita pasa por \texttt{locate}, que
exige que el párrafo exista y que la cita sea una subcadena literal de él. Cada rasgo se valida por
separado: una cita que falla hace repetir solo ese rasgo, una vez, y si vuelve a fallar queda como
fallido sin arrastrar a los demás. La partición en escenas debe cubrir todos los párrafos, en orden
y sin solapes, y la presencia de un personaje en una escena necesita una cita dentro de ella. Las
extracciones se guardan con la versión del \textit{prompt} y la huella del texto, y al relanzarlas
solo se reintentan los rasgos fallidos. \texttt{asg-audit-features} implementa la validación de la
Sección~\ref{sec:validacion-extractor}: escribe la hoja de revisión humana, con una fila por
ocurrencia, compara dos extracciones del mismo texto y contrasta los recuentos con el registro.

El estudio por pares (\texttt{study.py}) vive en una base SQLite. Al abrir la votación, el conjunto de
historias se congela con una huella SHA-256, y cada pregunta se calcula cuando el lector la pide, con
las reglas de la Sección~\ref{sec:conjunto}; la exportación es seudónima. Las fuerzas de
Bradley-Terry (\texttt{preferences.py}) se obtienen minimizando la verosimilitud con penalización de
cresta mediante el optimizador L-BFGS de SciPy~\cite{virtanen2020}, con 1.000 remuestreos de
lectores por defecto. El juez (\texttt{judge.py}) es la regresión logística de
scikit-learn~\cite{pedregosa2011}, con regularización $C = 1$ y sin término independiente; cada juez
guardado lleva la versión de su formato, y el ranking rechaza un juez de otra versión.
\texttt{asg-rank} lo aplica al corpus y prepara una ficha por cada una de las diez mejores, y
\texttt{asg-corpus}, de solo lectura, da el veredicto de elegibilidad de cada ejecución con su razón.

Para el análisis horizontal, \texttt{asg-matrix} genera la matriz en serie: antes de cada unidad
comprueba la cuota y, si no alcanza, para con un manifiesto al día; al relanzarla, salta las
unidades completas sin repetir llamadas. \texttt{asg-report-features} empareja cada celda con su
base, calcula las diferencias por rasgo y aplica la regla de ruido de la
Sección~\ref{sec:horizontal}.

\section{Interfaces}\label{sec:interfaces}

El bot de Telegram (\texttt{asg-telegram}), construido con python-telegram-bot~\cite{ptb2026}, corre
en la computadora del investigador. Para crear, ofrece cuatro opciones: petición libre o guiada paso
a paso, con el formato narrativo o con el simulado. Las peticiones entran en una cola SQLite con una
sola petición activa por persona, se genera una historia cada vez y cada una llega como archivo,
como texto y como audio, narrado con edge-tts. Un único módulo del bot importa el generador. Para
comparar hay dos flujos. Con \texttt{/evaluar}, cada persona compara a ciegas sus propias dos
historias, una de cada formato, con las tres preguntas; son votos de prueba que no entran en el
análisis. Con \texttt{/estudio} se vota el estudio de la Sección~\ref{sec:conjunto}: historias de
otras personas, solo como texto, sin título, encabezados ni audio, y la pregunta con A, B y «No
puedo decidir». Un candado por lector evita que una doble pulsación registre dos votos.

StageCraft (\texttt{asg-studio}) es una aplicación web local, servida con FastAPI~\cite{fastapi2026}
y escrita en HTML, CSS y módulos de JavaScript, sin paso de compilación. Sus vistas permiten crear
una obra viendo el \textit{prompt} exacto que se enviará, leer las ejecuciones guardadas con su
registro, también con los ojos de una visión, y comparar de dos a cuatro ejecuciones eje por eje,
indicando si la comparación está emparejada. La vista de \textbf{lectura ciega} muestra pares de
historias anónimas, X e Y, y guarda un voto por par; las respuestas de la aplicación no llevan
identificador, título ni configuración, y la clave de cada par se guarda aparte. Con ella se hicieron
las dos lecturas del 27 de octubre de 2026 que decidieron los brazos de escritura y de narración
(Sección~\ref{sec:opcionales}). StageCraft nunca escribe en las carpetas de las historias, solo
acepta escrituras con una cabecera propia y desde la propia máquina, y pinta como texto, nunca como
HTML, lo que produce el modelo.

\section{Pruebas y reproducibilidad}\label{sec:pruebas}

El 1 de noviembre de 2026, la suite completa pasó 907 pruebas y omitió las 2 que llaman a la API
real, que solo se ejecutan a petición porque consumen cuota; tardó algo más de dos minutos. Ninguna
otra prueba llama a un modelo: un proveedor falso recibe una secuencia de respuestas estructuradas,
puede forzar fallos en llamadas concretas y registra cada contexto que se le envía. Sobre una
ejecución completa con ese proveedor, varias pruebas comprueban los contratos del
Capítulo~\ref{chapter:generator}. \codigo{test_no_actor_ever_sees_the_plan_or_a_future_scene} reúne
los contextos de todos los actores y comprueba que en ninguno aparece el identificador ni el título
de un evento, ni el de un capítulo. \codigo{test_a_character_only_remembers_what_it_witnessed}
comprueba que cada registro de memoria procede de un turno que el personaje presenció.
\codigo{test_without_the_inventory_no_object_reaches_a_prompt_or_an_artifact} comprueba que, con el
inventario apagado, no se escribe ningún artefacto suyo ni llega a un \textit{prompt} ningún bloque
o esquema de objetos: un brazo apagado deja el generador como estaba. Y
\codigo{test_the_voice_only_changes_the_narration} genera dos ejecuciones que solo difieren en la
visión, el personaje y el tono, y exige que el plan, el guion, el reparto y la función sean idénticos
byte a byte y que las llamadas anteriores a la narración coincidan.

La reproducibilidad tiene un límite que conviene declarar: el modelo no es determinista, y dos
ejecuciones con el mismo \textit{prompt} y la misma versión producen historias distintas. Por eso la
evaluación trata ese ruido como un dato (Sección~\ref{sec:horizontal}) en lugar de suponer que una
ejecución se repite. Lo que sí se reproduce es todo lo que decide el código: dadas las mismas
respuestas del modelo, la validación, la percepción, la recuperación de la memoria, que no usa
\textit{embeddings} ni llamadas al modelo, y las medidas dan el mismo resultado. También se
reproduce el análisis: como cada ejecución conserva sus artefactos con su huella, la tabla de
rasgos, las preferencias, el juez y el ranking se recalculan desde el disco sin volver a generar, el
sorteo del estudio se repite con su semilla y cada historia declara la versión del generador, del
contrato y de los modelos que la produjeron.

Este capítulo deja descrito un sistema en el que las decisiones del Capítulo~\ref{chapter:generator}
tienen un lugar comprobable: los validadores y la memoria son funciones puras con pruebas, la
evaluación no puede importar el generador y cada ejecución queda en disco con su versión y su
huella. La cuota del servicio, administrada pero no evitada, acotó lo que se pudo generar. El
Capítulo~\ref{chapter:experiments} presenta lo que se midió con estas herramientas.
